Im Grenzfall besteht der Inhalt am Schluss vollständig aus Eis von $\vartheta_0 = \TnO$. Das Wasser (Masse $m_\text{W}$) kühlt auf $\vartheta_0$ ab und gefriert, das Eis (Masse $m_\text{E}$) erwärmt sich auf $\vartheta_0$. Mit $c_\text{W} = \cWO$ (FoTaBe S.~191), $c_\text{E} = \cEO$ und $L_\text{f} = \LfEO$ (FoTaBe S.~190) gilt
\begin{align}
 m_\text{E}\,c_\text{E}\left(\vartheta_0-\vartheta_\text{E}\right) &= m_\text{W}\left(c_\text{W}\left(\vartheta_\text{W}-\vartheta_0\right)+L_\text{f}\right)
\end{align}
Das Massenverhältnis ist
\begin{align}
 \frac{m_\text{E}}{m_\text{W}} &= \ratF \\
   &= \frac{\cW\cdot\left(\TWO-\TnO\right)+\LfE}{\cE\cdot\left(\TnO-\left(\TEO\right)\right)} \\
   &= \rat \approx \result{\ratP{0}{2}}
\end{align}
Man braucht also rund \ratP{0}{2}-mal mehr Eis als Wasser. Die Erstarrungswärme ist viel grösser als die Wärme, die das Eis bei einer Erwärmung um wenige Grad aufnehmen kann.
